%% kapa: exam=1 round=31 year=2020 period=2 form=A subject=appraisal_law subject_ko=감정평가관계법규 seq=29 number=29
%% source: Q-Net 감정평가사 기출문제 (cst00309) artlSeq=5207219 file=extracted/1008/5207219/A형/A형/2020년도 제31회 감정평가사 1차 시험_2교시_A형.pdf
%% conversion: kordoc 4.24.51 -> tools/gen1.py (auto); review=proofread-sonnet; flags=none
\begin{question}{29}
공간정보의 구축 및 관리 등에 관한 법령상 지적공부에 관한 내용으로 옳지 않은 것은?
\begin{choices}
  \item 지적소관청은 관할 시ㆍ도지사의 승인을 받은 경우 지적서고에 보존되어 있는 지적공부를 해당 청사 밖으로 반출할 수 있다.
  \item 지적공부를 정보처리시스템을 통하여 기록ㆍ저장한 경우 관할 시ㆍ도지사, 시장ㆍ군수 또는 구청장은 그 지적공부를 지적정보관리체계에 영구히 보존하여야 한다.
  \item 지적소관청은 부동산의 효율적 이용과 부동산과 관련된 정보의 종합적 관리ㆍ운영을 위하여 부동산종합공부를 관리ㆍ운영한다.
  \item 부동산종합공부를 열람하거나 부동산종합증명서를 발급받으려는 자는 지적소관청이나 읍ㆍ면ㆍ동의 장에게 신청할 수 있다.
  \item 지적전산자료를 신청하려는 자는 지적전산자료의 이용 또는 활용 목적 등에 관하여 미리 중앙지적위원회의 심사를 받아야 한다.
\end{choices}
\end{question}
\answer{5}
